\documentclass[11pt,letterpaper]{article}

% ---------- Idioma y codificación ----------
\usepackage[utf8]{inputenc}
\usepackage[T1]{fontenc}
\usepackage[spanish,es-tabla]{babel}

% ---------- Paquetes generales ----------
\usepackage[margin=2.5cm]{geometry}
\usepackage{xcolor}
\usepackage{enumitem}
\usepackage{fancyhdr}
\usepackage{titlesec}
\usepackage{tcolorbox}
\tcbuselibrary{skins,breakable}
\usepackage{listings}
\usepackage{graphicx}
\usepackage{array}
\usepackage{longtable}
\usepackage{hyperref}
\usepackage{xurl}
\usepackage{seqsplit}

% ---------- Colores del curso ----------
\definecolor{brandblue}{HTML}{1F4E78}
\definecolor{brandgreen}{HTML}{1DB954}
\definecolor{lightgray}{HTML}{F4F6F8}
\definecolor{codebg}{HTML}{F5F5F5}

% ---------- Encabezado y pie de página ----------
\pagestyle{fancy}
\fancyhf{}
\renewcommand{\headrulewidth}{0.4pt}
\fancyhead[L]{\small Laboratorio de Simulación 2 --- Analítica de Datos}
\fancyhead[R]{\small Módulo 2: Transformación y visualización de datos}
\fancyfoot[C]{\small\thepage}

% ---------- Estilo de secciones ----------
\titleformat{\section}
  {\normalfont\Large\bfseries\color{brandblue}}
  {\thesection}{0.6em}{}
\titleformat{\subsection}
  {\normalfont\large\bfseries\color{brandblue}}
  {\thesubsection}{0.6em}{}

% ---------- Cajas de actividades y notas ----------
\newtcolorbox{actividad}[1][]{
  breakable, enhanced, colback=lightgray, colframe=brandblue,
  boxrule=0.8pt, arc=2pt, left=6pt, right=6pt, top=5pt, bottom=5pt,
  title={\faPencil\ Actividad}, fonttitle=\bfseries\color{white},
  colbacktitle=brandblue, #1
}
\newtcolorbox{nota}[1][]{
  breakable, enhanced, colback=white, colframe=brandgreen,
  boxrule=0.8pt, arc=2pt, left=6pt, right=6pt, top=5pt, bottom=5pt,
  title={Nota}, fonttitle=\bfseries\color{white},
  colbacktitle=brandgreen, #1
}

% ---------- Código / comandos ----------
\lstset{
  basicstyle=\ttfamily\small,
  backgroundcolor=\color{codebg},
  breaklines=true,
  frame=single,
  framerule=0.3pt,
  rulecolor=\color{gray!50},
  columns=fullflexible,
  keepspaces=true,
}

% ---------- Metadatos del PDF ----------
\hypersetup{
  pdftitle={Laboratorio de Simulación 2 - Analítica de Datos},
  pdfauthor={Diego Fabián Collazos Huertas},
  pdfsubject={Curso Analítica de Datos - Módulo 2},
  colorlinks=true,
  linkcolor=brandblue,
  urlcolor=brandblue,
  citecolor=brandblue,
}

% Sustituto simple para el icono de lápiz (evita depender de fontawesome)
\newcommand{\faPencil}{\raisebox{-0.1em}{\textbf{\Large\textbullet}}}

\title{\vspace{-2cm}}
\date{}

\begin{document}

% ============================================================
% PORTADA
% ============================================================
\begin{titlepage}
\centering
\vspace*{1.5cm}

{\Large Universidad Nacional de Colombia --- Sede Manizales}\\[0.2cm]
{\large Facultad de Ingeniería y Arquitectura}\\[0.2cm]
{\large Departamento de Ingeniería Eléctrica, Electrónica y Computación}\\[1.2cm]

\rule{\linewidth}{1pt}\\[0.5cm]
{\huge\bfseries\color{brandblue} Laboratorio de Simulación 2}\\[0.4cm]
{\Large Implementación en Orange Data Mining}\\[0.5cm]
\rule{\linewidth}{1pt}\\[1cm]

{\large \textbf{Curso:} Analítica de Datos (Aprendizaje de máquina) --- Código 4200729}\\[0.3cm]
{\large \textbf{Módulo:} 2 --- Transformación y visualización de datos}\\[0.3cm]
{\large \textbf{Docente:} Diego Fabián Collazos Huertas}\\[0.3cm]
{\large \textbf{Modalidad:} trabajo en parejas}\\[1.5cm]

\vfill

{\small Este laboratorio hace parte de la nota de \textbf{Laboratorios de simulación (45\% de la evaluación)}, según el programa de la asignatura (\texttt{PROGRAMA.md}).}

\end{titlepage}

% ============================================================
% OBJETIVOS
% ============================================================
\section{Objetivos}

\begin{itemize}[leftmargin=1.2cm]
  \item Reproducir en \textbf{Orange Data Mining} --- mediante programación visual, sin escribir código --- los mismos análisis que se implementaron en Python en los notebooks \texttt{\seqsplit{2.2\_Preprocesamiento\_Transformacion\_Ejemplos.ipynb}} y \texttt{\seqsplit{2.4\_Visualizacion\_Datos\_Ejemplos.ipynb}}.
  \item Diagnosticar datos faltantes y aplicar el tratamiento adecuado (excluir, imputar) usando los widgets de preprocesamiento de Orange.
  \item Dividir un dataset en entrenamiento/prueba y aplicar la \textbf{misma} transformación aprendida en entrenamiento sobre la prueba, evitando la fuga de información.
  \item Discretizar una variable numérica de las tres formas vistas en clase (ancho uniforme, frecuencia uniforme, supervisada) y codificar variables categóricas.
  \item Elegir y construir las visualizaciones adecuadas (histograma, diagrama de caja, barras, dispersión, serie temporal) para responder preguntas sobre un dataset real.
\end{itemize}

\begin{nota}
Este taller \textbf{no reemplaza} los notebooks de Python; los complementa. La idea es que compares, para el mismo problema, qué tan distinto es resolverlo con código (Pandas/scikit-learn/Matplotlib) frente a resolverlo con programación visual (Orange) --- ver \texttt{1.2\_Orange\_Data\_Mining.md} para instalar la herramienta si aún no lo has hecho.
\end{nota}

% ============================================================
% INSTRUCCIONES GENERALES
% ============================================================
\section{Instrucciones generales}

\begin{enumerate}[leftmargin=1.2cm]
  \item \textbf{Este laboratorio se desarrolla en parejas} (2 estudiantes), asignadas o conformadas según indique el docente.
  \item Para cada parte del laboratorio, arma el flujo (\textit{workflow}) solicitado en el lienzo de Orange.
  \item Cada vez que veas el ícono \faPencil\ \textbf{Actividad}, documenta la respuesta \textbf{directamente sobre el lienzo de Orange}, usando la herramienta de texto (clic derecho sobre un espacio vacío del lienzo $\to$ \textit{Text}, o el ícono ``T'' de la barra de herramientas) para agregar un cuadro de texto junto a los widgets correspondientes. No se entregan respuestas en un documento aparte: el workflow comentado \textbf{es} la respuesta.
  \item Además, toma una captura de pantalla del resultado de cada actividad (el widget correspondiente ya abierto, mostrando el gráfico o la tabla) para incluirla en el PDF de evidencias.
  \item Consulta la sección \textbf{Entregables} (al final de este documento) para el detalle de qué debes entregar y cómo.
\end{enumerate}

\begin{nota}
Algunos widgets de este laboratorio (\textit{Group By}, \textit{Line Plot}) son relativamente recientes en Orange. Si no los encuentras en tu barra de widgets, actualiza Orange a la versión más reciente (\textit{Help $\to$ Check for Updates}, o reinstalando desde \url{https://orangedatamining.com/download/}).
\end{nota}

% ============================================================
% PARTE 1
% ============================================================
\section{Parte 1 --- Preprocesamiento y transformación en Orange}

Esta parte reproduce en Orange lo desarrollado en el notebook \texttt{\seqsplit{2.2\_Preprocesamiento\_Transformacion\_Ejemplos.ipynb}}, usando el mismo dataset: el \textbf{Titanic}.

\subsection{1.1 Cargar el dataset y diagnosticar datos faltantes}

\begin{enumerate}[leftmargin=1.2cm]
  \item Reutiliza el archivo \texttt{.csv} del Titanic que ya descargaste con \texttt{kagglehub} al ejecutar el notebook \texttt{2.2} (queda guardado en la ruta que imprime la celda de descarga), o descárgalo directamente desde \url{https://www.kaggle.com/datasets/yasserh/titanic-dataset}.
  \item Abre Orange y crea un workflow nuevo. Arrastra el widget \textbf{File} al lienzo y carga el archivo.
  \item En el editor de columnas del widget \textit{File}, marca la columna \texttt{Survived} con el rol \textbf{target} (clase) --- la necesitaremos más adelante para la discretización supervisada.
  \item Agrega el widget \textbf{Feature Statistics} (categoría \textit{Data}) y conéctalo a \textit{File}. Esta tabla muestra, entre otras cosas, el porcentaje de valores faltantes de cada columna.
\end{enumerate}

\begin{actividad}
Compara la columna de faltantes de \textit{Feature Statistics} con el diagnóstico del notebook (\texttt{Cabin} $\approx$ 77\%, \texttt{Age} $\approx$ 20\%, \texttt{Embarked} $\approx$ 0.2\%). En un cuadro de texto junto al widget, escribe si los porcentajes coinciden. Incluye en tu PDF de evidencias una captura de \textit{Feature Statistics}.
\end{actividad}

\subsection{1.2 Dividir train/test y tratar los datos sin fuga de información}

Como en el notebook (sección 3), \textbf{dividimos antes de calcular cualquier estadístico de imputación}, para que ningún dato de prueba influya en los valores usados para imputar.

\begin{enumerate}[leftmargin=1.2cm]
  \item Agrega el widget \textbf{Select Columns} conectado a \textit{File}, y mueve la columna \texttt{Cabin} al grupo \textit{Excluded} (equivalente a \texttt{drop(columns=['Cabin'])}).
  \item Agrega el widget \textbf{Data Sampler}, conéctalo a \textit{Select Columns}. Configúralo con \textbf{Fixed proportion of data: 80\%} y activa la casilla \textbf{Stratify sample}. Esto reparte los datos en dos salidas: \textit{Data Sample} (entrenamiento) y \textit{Remaining Data} (prueba) --- equivalente a \texttt{train\_test\_split(..., stratify=titanic['Survived'])}.
  \item Agrega un primer widget \textbf{Preprocess} y conéctale la salida \textit{Data Sample} (entrenamiento). Configúralo con \textbf{Impute Missing Values}: \texttt{Age} $\to$ \textit{Median}, y las variables categóricas (incluida \texttt{Embarked}) $\to$ \textit{Most frequent value}.
  \item Este primer \textit{Preprocess} tiene dos salidas: \textbf{Preprocessed Data} (el entrenamiento ya tratado) y \textbf{Preprocessor} (la ``receta'' aprendida: la mediana de \texttt{Age} y la moda de \texttt{Embarked}, calculadas \textbf{solo} con el entrenamiento).
  \item Agrega un \textbf{segundo} widget \textbf{Preprocess}. Conéctale la salida \textit{Remaining Data} (prueba) de \textit{Data Sampler} como \textbf{Data}, y la salida \textit{Preprocessor} del primer \textit{Preprocess} como \textbf{Preprocessor}. Este segundo widget \textbf{no vuelve a calcular nada}: solo aplica al conjunto de prueba la misma mediana y moda aprendidas en entrenamiento.
\end{enumerate}

\begin{nota}
Este par de widgets \textit{Preprocess} conectados entre sí es el equivalente visual exacto de ``ajustar (\textit{fit}) en entrenamiento y transformar (\textit{transform}) en prueba con los mismos parámetros'', que en el notebook se resuelve calculando \texttt{mediana\_age} y \texttt{moda\_embarked} solo con \texttt{train} y aplicándolas también a \texttt{test}.
\end{nota}

\begin{actividad}
Conecta un widget \textbf{Data Table} a cada una de las dos salidas \textit{Preprocessed Data} (la del entrenamiento y la de la prueba) y verifica que ninguna de las dos columnas tratadas (\texttt{Age}, \texttt{Embarked}) tenga valores faltantes. En un cuadro de texto junto al segundo \textit{Preprocess}, anota el valor con el que quedó imputada \texttt{Age} y compáralo con la mediana que calculó el notebook ($\approx$ 28.5 años). Incluye en tu PDF de evidencias una captura del workflow completo de esta sección (\textit{Select Columns} $\to$ \textit{Data Sampler} $\to$ los dos \textit{Preprocess} $\to$ las dos \textit{Data Table}).
\end{actividad}

\subsection{1.3 Discretización de la edad y codificación de categóricas}

\begin{enumerate}[leftmargin=1.2cm]
  \item Agrega el widget \textbf{Discretize} (categoría \textit{Data}) conectado a la salida \textit{Preprocessed Data} del entrenamiento. Selecciona la variable \texttt{Age} y prueba, uno a la vez, los tres métodos de discretización disponibles:
  \begin{itemize}
    \item \textbf{Equal frequency} (frecuencia uniforme) --- equivalente a \texttt{pd.qcut} en el notebook.
    \item \textbf{Equal width} (ancho uniforme) --- equivalente a \texttt{pd.cut} con bordes fijos.
    \item \textbf{Entropy-MDL} --- discretización \textbf{supervisada}: usa la variable objetivo (\texttt{Survived}, que ya marcaste como \textit{target} en el punto 1.1) para elegir el corte que mejor separa las clases, igual que la búsqueda por entropía del notebook.
  \end{itemize}
  \item Agrega un widget \textbf{Continuize} (categoría \textit{Data}) conectado también a la salida del entrenamiento. Selecciona \texttt{Embarked} y prueba las opciones \textbf{One feature per value} (equivalente a \textit{one-hot}) y \textbf{One feature per value, ignore the first value} (equivalente a \textit{dummy}, la codificación $N-1$).
\end{enumerate}

\begin{actividad}
En un cuadro de texto junto al widget \textit{Discretize}, anota el corte que encontró el método \textbf{Entropy-MDL} para \texttt{Age} y compáralo con el umbral que encontró el notebook ($\approx$ 5--6.5 años, el patrón de ``mujeres y niños primero''). Junto al widget \textit{Continuize}, anota cuántas columnas nuevas genera cada opción para \texttt{Embarked} y compáralas con las columnas \texttt{Embarked\_C/Q/S} (one-hot) y \texttt{Embarked\_Q/S} (dummy) del notebook. Incluye en tu PDF de evidencias una captura de cada widget con el método/opción que mejor haya reproducido el resultado del notebook.
\end{actividad}

% ============================================================
% PARTE 2
% ============================================================
\section{Parte 2 --- Visualización en Orange}

Esta parte reproduce en Orange lo desarrollado en el notebook \texttt{\seqsplit{2.4\_Visualizacion\_Datos\_Ejemplos.ipynb}}, usando el mismo dataset: \textbf{Video Game Sales}.

\subsection{2.1 Cargar el dataset y explorar la distribución}

\begin{enumerate}[leftmargin=1.2cm]
  \item Reutiliza el archivo \texttt{.csv} de Video Game Sales que ya descargaste con \texttt{kagglehub} al ejecutar el notebook \texttt{2.4}, o descárgalo directamente desde \url{https://www.kaggle.com/datasets/gregorut/videogamesales}.
  \item Carga el archivo en un nuevo workflow con el widget \textbf{File}.
  \item Agrega el widget \textbf{Distributions}, conéctalo a \textit{File}, y selecciona la variable \texttt{Global\_Sales}. Prueba distintos valores del control de número de intervalos (\textit{Bin width} / número de \textit{bins}), igual que en el notebook.
  \item Agrega el widget \textbf{Feature Statistics} conectado a \textit{File}: muestra directamente la \textbf{media}, la \textbf{mediana} y la \textbf{asimetría (skewness)} de \texttt{Global\_Sales}.
\end{enumerate}

\begin{actividad}
Compara el histograma de \textit{Distributions} con los del notebook (¿también se ve una cola larga hacia la derecha con pocos \textit{bins} y se nota mejor con más \textit{bins}?). En un cuadro de texto junto a \textit{Feature Statistics}, compara los valores de media, mediana y asimetría con los que calculó el notebook (media $\approx 0.54$, mediana $\approx 0.17$, asimetría muy alta y positiva) y explica en una frase qué significa que la media sea mucho mayor que la mediana. Incluye en tu PDF de evidencias una captura de cada widget.
\end{actividad}

\subsection{2.2 Diagrama de caja, valores atípicos y variables categóricas}

\begin{enumerate}[leftmargin=1.2cm]
  \item Agrega el widget \textbf{Box Plot} conectado a \textit{File}, y selecciona \texttt{Global\_Sales}. Orange dibuja la caja, la mediana y marca los valores atípicos según la misma regla del IQR usada en el notebook.
  \item Agrega el widget \textbf{Bar Plot} conectado a \textit{File}, y selecciona \texttt{Genre} para ver el conteo de juegos por género.
\end{enumerate}

\begin{actividad}
En un cuadro de texto junto al \textit{Box Plot}, anota cuántos valores atípicos identifica Orange y compáralos con el conteo de la regla del IQR del notebook ($\approx$ 11\% de los juegos). Junto al \textit{Bar Plot}, indica si el género más frecuente coincide con el del notebook (\textit{Action}). Incluye en tu PDF de evidencias una captura de cada widget.
\end{actividad}

\subsection{2.3 Relaciones entre variables y serie temporal}

\begin{enumerate}[leftmargin=1.2cm]
  \item Agrega el widget \textbf{Scatter Plot} conectado a \textit{File}: en el eje X selecciona \texttt{NA\_Sales}, en el eje Y selecciona \texttt{EU\_Sales}.
  \item Agrega el widget \textbf{Correlations} conectado a \textit{File}: ordena automáticamente todos los pares de variables numéricas según su coeficiente de correlación --- es el equivalente a leer el mapa de calor de correlaciones del notebook, pero en forma de tabla.
  \item Agrega el widget \textbf{Group By} (categoría \textit{Transform}) conectado a \textit{File}: agrupa por \texttt{Year} y agrega \texttt{Global\_Sales} con la función \textit{Sum} --- es el equivalente visual de \texttt{\seqsplit{groupby('Year')['Global\_Sales'].sum()}} en el notebook.
  \item Conecta la salida de \textit{Group By} a un widget \textbf{Line Plot} para ver la serie temporal de ventas globales por año.
\end{enumerate}

\begin{actividad}
En un cuadro de texto junto al \textit{Scatter Plot}, describe la dirección y forma de la relación entre \texttt{NA\_Sales} y \texttt{EU\_Sales}, y compárala con la correlación que reporta el widget \textit{Correlations} (el notebook encontró $\approx 0.77$). Junto al \textit{Line Plot}, indica si el año de mayores ventas coincide con el que encontró el notebook (2008). Incluye en tu PDF de evidencias una captura de los cuatro widgets de esta sección.
\end{actividad}

% ============================================================
% ENTREGABLES
% ============================================================
\section{Entregables}

La pareja debe subir al repositorio del curso (a través de la página de entregas) \textbf{dos tipos de archivo}:

\begin{enumerate}[leftmargin=1.2cm]
  \item \textbf{Los archivos \texttt{.ows} de los workflows armados en Orange} (uno por cada parte del laboratorio), \textbf{comentados con la herramienta de texto de Orange} tal como se indicó en cada actividad. El workflow comentado --- no un documento aparte --- es donde se responden las preguntas de cada actividad.
  \item \textbf{Un único PDF con las capturas de pantalla} de los resultados obtenidos en cada una de las 6 actividades (\faPencil), claramente identificadas (por ejemplo, ``Actividad 1 --- Feature Statistics del Titanic'').
\end{enumerate}

Nombra los archivos con los apellidos de ambos integrantes de la pareja, por ejemplo: \texttt{\seqsplit{apellido1\_apellido2\_parte1.ows}}, \texttt{\seqsplit{apellido1\_apellido2\_parte2.ows}} y \texttt{\seqsplit{apellido1\_apellido2\_evidencias.pdf}}. Sube estos archivos desde la página de entregas del curso, seleccionando la categoría \textbf{Laboratorio de Simulación 2}.

\begin{actividad}[title={\faPencil\ Reflexión final}]
Como cuadro de texto de cierre en el workflow de la \textbf{Parte 1}, escriban en pareja una breve reflexión (5-8 líneas): ¿qué fue más fácil de entender en Orange que en Python (o viceversa) al dividir train/test y evitar la fuga de información? ¿Y al discretizar y codificar variables?
\end{actividad}

\begin{nota}
Este laboratorio corresponde a la segunda entrega del componente \textbf{Laboratorios de simulación (45\% de la nota)}. Revisa la fecha de entrega vigente en \texttt{entregas/README.md}.
\end{nota}

\end{document}
